\chapter{Introducción}
\label{ch:introduccion}

Este informe documenta el desarrollo y la resolución del Ejercicio Práctico N° 2 del Trabajo Práctico N° 2.
Las actividades propuestas en esta etapa tienen como objetivo principal
estudiar y comprender la configuración y puesta en marcha del convertidor analógico-digital (ADC1) en el
microcontrolador STM32F103C8T6 bajo un esquema de programación Bare Metal puro.

El propósito general de este ejercicio es implementar la adquisición de señales analógicas y su posterior
procesamiento digital. Para ello, se busca comprender el funcionamiento del periférico de conversión, la
configuración de canales múltiples, el establecimiento de tiempos de muestreo y factores de escala, y la
representación de las mediciones obtenidas, la cual se logra mediante el uso de
5 LEDs. A su vez, se integra el control de eventos por hardware
para alternar dinámicamente entre las fuentes de señal de entrada, vinculando las técnicas de interrupciones
externas estudiadas en el ejercicio anterior con los métodos de muestreo de este periférico.
